\subsection{O-013: the Hessian family at every outer degree}

Let $n\geq2$, $1\leq r\leq\min(d,n^2)$, and put
\[
          \lambda_{d,r}=(nd-2(r-1),2^{r-1})\vdash nd.
\]
Then
\[
           K_n(\lambda_{d,r})=a_{\lambda_{d,r}}(d[n])=1.
                                                        \tag{O-013.1}
\]
For internal padding and $n>m\geq2$, one moreover has
\[
 q^m_{\lambda_{d,r}}(d[n])=
 \begin{cases}1&r\leq m^2,\\0&r>m^2.\end{cases}          \tag{O-013.2}
\]
There is no upper bound on the outer degree $d$. Thus this
extends O-011 beyond the maximal-row case $d=r$ and gives
genuine positive-permanent examples even at degrees beyond
the independent quadratic Chow bound M-010.

\begin{proof}
The one-row case is the power of the first coefficient, so
assume $r\geq2$. Apply F-014 with second row $2$:
\[
 a_{\lambda_{d,r}}(d[n])=
 a_{\mu}(d[2]),\qquad
       \mu=(2(d-r+1),2^{r-1}).
                                                        \tag{O-013.3}
\]
We give the short multiplicity-free argument needed on the
right, rather than presuming a formula for all quadratic
plethysms. In the space of quadratic forms on $\mathbb C^r$,
the upper triangular group has a dense open orbit. Indeed,
on the open set where all leading principal minors of the
symmetric coefficient matrix are nonzero, symmetric Gaussian
elimination gives an upper-triangular congruence to a diagonal
matrix with nonzero entries. Over $\mathbb C$ a further
diagonal congruence changes that matrix to the identity.
Consequently the ratio of two nonzero polynomial highest-weight
functions of the same weight is a rational invariant of the
upper triangular group, and is constant on its open orbit.
The two functions are proportional. Every irreducible
multiplicity in the coordinate ring of quadratic forms is
therefore at most one.

For a degree-$n$ form $F$, let $W_r(F)$ be the Hessian minor
used in O-011, and define
\[
                  H_{d,r}(F)=W_r(F)F(e_1)^{d-r}.
                                                        \tag{O-013.4}
\]
This is a degree-$d$ polynomial function. The chain-rule
calculation in O-011 and the identity
$F(ge_1)=g_{11}^{n}F(e_1)$ for upper triangular $g$ show
that it is a highest-weight vector of weight
\[
  (r(n-2)+2+n(d-r),2^{r-1})=\lambda_{d,r}.
\]
For $n=2$, evaluation on a nonsingular diagonal quadratic
form with nonzero first coefficient makes (O-013.4) nonzero.
Thus the multiplicity at $\mu$ in (O-013.3) is exactly one,
and so is the ambient multiplicity at $\lambda_{d,r}$.

At the determinant point, O-011 supplies for every $r\leq n^2$
a change of coordinates under which $W_r$ is nonzero and
$F(e_1)$ is nonzero: the first direction corresponds to the
identity matrix, and the other directions span a nondegenerate
Hessian subspace containing it. Hence (O-013.4) is nonzero
on the determinant orbit, proving $K_n\geq1$. The ambient
upper bound gives (O-013.1).

For the internally padded permanent, O-011 supplies the same
two nonvanishing conditions for every $r\leq m^2$, including
the separately checked exceptional pair $(m,n)=(3,4)$.
Thus (O-013.4) is nonzero on that orbit as well, and
$q\geq1$. Again the ambient upper bound gives $q=1$.
If $r>m^2$, the permanent depends on only $m^2$ variables,
so the length restriction gives $q=0$.
\end{proof}

This is an exact scalar family at all outer degrees, with
an explicit surviving highest-weight function. It does not
claim that occurrence determines multiplicity when the ambient
coefficient exceeds one. The determinant-side proof uses
the actual orbit, so no replacement of closure multiplicities
by symmetric or ordinary rectangular coefficients occurs.

\paragraph{Geometric dictionary.}
Quadratic forms are points of the affine scheme of symmetric
bilinear matrices. Triangular congruence gives the displayed
open orbit after base change to $\mathbb C$; the resulting
character semi-invariants are torus-weight sections in the
grading/comodule correspondence of SGA~3, I, Proposition~4.7.3,
pp.~35--36. The Hessian and evaluation polynomials define
sections of the indicated Schur-type summands in the coordinate
algebra of forms. Their functorial decompositions are Macdonald,
I Appendix~A, (5.3)--(5.6), pp.~155--156, and the locally free
$K_0$ framework is SGA~6, VI, Theorem~3.3, pp.~392--393.
The scheme maps and derivative polynomials have integral
coefficients, but the Borel orbit argument, complete
reducibility, and Hessian nondegeneracy used here are
characteristic-zero assertions. Nonzero restriction to the
actual orbit means that the corresponding section survives
the coherent ideal quotient of its closure.
